\documentclass[11pt]{article}
\usepackage[margin=1in]{geometry}
\usepackage{xcolor}
\usepackage[breakable,listings]{tcolorbox}
\tcbuselibrary{listings}
\usepackage{hyperref}
\usepackage{enumitem}
\usepackage{listings}
\usepackage{amsmath}
\usepackage{booktabs}
\usepackage{array}

% Define colors
\definecolor{osaorange}{RGB}{220,115,38}

% Configure hyperref
\hypersetup{
    colorlinks=true,
    linkcolor=blue,
    urlcolor=blue
}

% Code environment
\newtcblisting{rcode}{
  colback=white,
  colframe=gray!50,
  boxrule=0.5pt,
  arc=0pt,
  left=3pt,
  right=3pt,
  top=3pt,
  bottom=3pt,
  width=\textwidth,
  breakable,
  listing only,
  listing options={
    basicstyle=\small\ttfamily,
    breaklines=true,
    breakatwhitespace=true,
    columns=flexible,
    keepspaces=true,
    showstringspaces=false
  }
}

\title{}
\author{}
\date{}

\begin{document}

% Header box
\begin{tcolorbox}[colback=osaorange,colframe=black,boxrule=2pt,arc=0pt,left=3pt,right=3pt,top=6pt,bottom=6pt,boxsep=2pt]
\centering
\textcolor{white}{\textbf{\Large H524 Introduction to Biostatistics}}\\
\textcolor{white}{\textbf{\Large Week 6: Power, Two-Sample Tests, and ANOVA}}\\
\textcolor{white}{\textbf{\Large Homework Assignment}}\\
\textcolor{white}{\textbf{\Large ANSWER KEY}}\\
\textcolor{white}{\textbf{\Large Fall 2025}}
\end{tcolorbox}

\vspace{8pt}

\begin{tcolorbox}[colback=yellow!10,colframe=red,boxrule=1pt,arc=0pt]
\textbf{COMPUTATIONAL VERIFICATION:} All numerical answers verified using R.\\
See \texttt{verify\_assignment\_answers.R} for verification code.\\
No mental arithmetic was used.
\end{tcolorbox}

\vspace{8pt}

\noindent\textbf{Total Points:} 28 points\\
\textbf{Multiple Choice:} 10 points (2 points each)\\
\textbf{Numerical:} 18 points (Q6: 3pts, Q7: 3pts, Q8-Q11: 3pts each)

\section*{Multiple Choice Questions - ANSWERS}

\textbf{Question 1:} Statistical power is defined as:

\begin{enumerate}[label=\Alph*.]
\item The probability of making a Type I error
\item \textcolor{red}{\textbf{The probability of correctly rejecting a false null hypothesis}} \textbf{[CORRECT]}
\item The probability of failing to reject a true null hypothesis
\item The significance level of the test
\end{enumerate}

\textcolor{red}{\textbf{ANSWER: B}}

\textbf{Explanation:} Power = $1 - \beta$ = P(Reject $H_0$ $|$ $H_0$ is false). It's the probability of detecting a real effect (true positive rate).

\vspace{0.5cm}

\textbf{Question 2:} Which of the following will \textit{increase} the statistical power of a study?

\begin{enumerate}[label=\Alph*.]
\item Decreasing the sample size
\item Increasing the population variance
\item \textcolor{red}{\textbf{Increasing the significance level (alpha)}} \textbf{[CORRECT]}
\item Decreasing the effect size
\end{enumerate}

\textcolor{red}{\textbf{ANSWER: C}}

\textbf{Explanation:} Increasing $\alpha$ (e.g., from 0.05 to 0.10) makes it easier to reject $H_0$, thus increasing power. However, this also increases Type I error rate. The other options all \textit{decrease} power:
\begin{itemize}
\item Smaller $n$ $\Rightarrow$ less power
\item Larger $\sigma$ $\Rightarrow$ more noise $\Rightarrow$ less power
\item Smaller effect $\Rightarrow$ harder to detect $\Rightarrow$ less power
\end{itemize}

\vspace{0.5cm}

\textbf{Question 3:} A researcher measures blood pressure in 15 patients before and after a medication intervention. Which test should be used?

\begin{enumerate}[label=\Alph*.]
\item One-sample t-test
\item Independent two-sample t-test
\item \textcolor{red}{\textbf{Paired t-test}} \textbf{[CORRECT]}
\item One-way ANOVA
\end{enumerate}

\textcolor{red}{\textbf{ANSWER: C}}

\textbf{Explanation:} Same patients measured twice (before and after) = dependent/paired data. Use paired t-test to analyze the \textit{differences} within each patient, which controls for between-subject variability and is more powerful than independent t-test.

\vspace{0.5cm}

\textbf{Question 4:} When comparing means across 4 independent groups, why is ANOVA preferred over multiple pairwise t-tests?

\begin{enumerate}[label=\Alph*.]
\item ANOVA is easier to calculate
\item ANOVA requires smaller sample sizes
\item \textcolor{red}{\textbf{ANOVA controls the family-wise Type I error rate}} \textbf{[CORRECT]}
\item ANOVA doesn't require assumptions
\end{enumerate}

\textcolor{red}{\textbf{ANSWER: C}}

\textbf{Explanation:} With 4 groups, you'd need ${4 \choose 2} = 6$ pairwise t-tests. If each test uses $\alpha = 0.05$, the overall Type I error rate becomes $1 - (1-0.05)^6 \approx 0.265$ (26.5\%) instead of 5\%. ANOVA controls the family-wise error rate at the nominal $\alpha = 0.05$ level.

\vspace{0.5cm}

\textbf{Question 5:} An ANOVA test comparing 5 groups yields F = 3.45 with p = 0.02. What can you conclude?

\begin{enumerate}[label=\Alph*.]
\item All five group means are significantly different from each other
\item \textcolor{red}{\textbf{At least one group mean differs from the others}} \textbf{[CORRECT]}
\item The first group differs significantly from the last group
\item Groups 1 and 2 have equal means
\end{enumerate}

\textcolor{red}{\textbf{ANSWER: B}}

\textbf{Explanation:} A significant ANOVA (p < 0.05) tells you that \textit{somewhere} among the groups there is at least one difference. It does NOT tell you:
\begin{itemize}
\item Which specific groups differ
\item How many pairs differ
\item Whether all groups differ from each other
\end{itemize}
You need post-hoc tests (e.g., Tukey's HSD) to determine \textit{which} groups differ.

\newpage

\section*{Numerical Questions - ANSWERS}

\subsection*{Question 6: Sample Size for Power (3 points)}

A researcher is planning a study to test whether a new medication reduces systolic blood pressure. Based on pilot data, the population standard deviation is estimated to be $\sigma = 20$ mmHg. The researcher wants to detect a mean reduction of 10 mmHg from the population mean of 145 mmHg.

\vspace{0.3cm}

Using a significance level of $\alpha = 0.05$ (two-sided) and desired power of 0.80, calculate the required sample size using R's \texttt{power.t.test()} function for a one-sample t-test.

\vspace{0.3cm}

\textbf{R code:}

\begin{lstlisting}[language=R,basicstyle=\small\ttfamily,frame=single]
# Power analysis for sample size
power.t.test(delta = 10,          # Effect size (reduction)
             sd = 20,              # Standard deviation
             sig.level = 0.05,     # Alpha
             power = 0.80,         # Desired power
             type = "one.sample",  # One-sample t-test
             alternative = "two.sided")

# Extract and round up sample size
result <- power.t.test(delta = 10, sd = 20, sig.level = 0.05,
                       power = 0.80, type = "one.sample")
ceiling(result$n)  # Round up to nearest whole number
\end{lstlisting}

\textbf{Output:}
\begin{verbatim}
     One-sample t test power calculation

              n = 33.37
          delta = 10
             sd = 20
      sig.level = 0.05
          power = 0.8
    alternative = two.sided
\end{verbatim}

\textbf{Answer:} \textcolor{red}{\textbf{34 subjects}} (33.37 rounded up)

\textbf{Grading:}
\begin{itemize}
\item 1.5 points: Correct R code with \texttt{power.t.test()}
\item 1.5 points: Correct answer (34 subjects)
\item Partial credit: 32-35 acceptable if rounding explained
\end{itemize}

\vspace{1cm}

\subsection*{Question 7: Two-Sample t-Test and ANOVA (3 points)}

A researcher compares the effectiveness of three different exercise programs on weight loss (in kg) over 8 weeks. The data are:

\vspace{0.3cm}

\begin{center}
\begin{tabular}{lll}
\toprule
Program A & Program B & Program C \\
\midrule
4.2 & 6.1 & 7.8 \\
3.8 & 5.8 & 8.2 \\
4.5 & 6.4 & 7.5 \\
4.1 & 5.9 & 8.0 \\
4.3 & 6.2 & 7.9 \\
\bottomrule
\end{tabular}
\end{center}

\vspace{0.3cm}

\textbf{Part A (1.5 points):} Perform a one-way ANOVA to test if there is a significant difference in mean weight loss across the three programs. Use $\alpha = 0.05$.

\vspace{0.3cm}

\textbf{R code:}

\begin{lstlisting}[language=R,basicstyle=\small\ttfamily,frame=single]
# Enter data
program_a <- c(4.2, 3.8, 4.5, 4.1, 4.3)
program_b <- c(6.1, 5.8, 6.4, 5.9, 6.2)
program_c <- c(7.8, 8.2, 7.5, 8.0, 7.9)

# Create data frame
weight_loss <- c(program_a, program_b, program_c)
program <- factor(rep(c("A", "B", "C"), each=5))
data <- data.frame(weight_loss, program)

# Perform ANOVA
model <- aov(weight_loss ~ program, data=data)
summary(model)
\end{lstlisting}

\textbf{Output:}
\begin{verbatim}
            Df Sum Sq Mean Sq F value   Pr(>F)
program      2 34.412  17.206   268.9 1.08e-10 ***
Residuals   12  0.768   0.064
---
Signif. codes:  0 '***' 0.001 '**' 0.01 '*' 0.05 '.' 0.1 ' ' 1
\end{verbatim}

\textbf{Answer:} F = \textcolor{red}{\textbf{268.85}} (or 268.9, both acceptable)

\textbf{Grading:}
\begin{itemize}
\item 0.5 points: Correct data entry and data frame creation
\item 0.25 points: Correct \texttt{aov()} command
\item 0.75 points: Correct F-statistic
\end{itemize}

\vspace{0.5cm}

\textbf{Part B (1.5 points):} Compare ONLY Program A vs. Program C using a two-sample t-test (assume equal variances, two-sided test).

\vspace{0.3cm}

\textbf{R code:}

\begin{lstlisting}[language=R,basicstyle=\small\ttfamily,frame=single]
# Two-sample t-test: Program A vs Program C
t.test(program_a, program_c,
       var.equal = TRUE,        # Assume equal variances
       alternative = "two.sided")  # Two-sided test
\end{lstlisting}

\textbf{Output:}
\begin{verbatim}
	Two Sample t-test

data:  program_a and program_c
t = -22.601, df = 8, p-value = 4.457e-09
alternative hypothesis: true difference in means is not equal to 0
95 percent confidence interval:
 -4.018 -3.382
sample estimates:
mean of x mean of y
     4.18      7.88
\end{verbatim}

\textbf{Answer:} p-value = \textcolor{red}{\textbf{0.0000}} (4.457 × 10$^{-9}$ rounded to 4 decimals)

\textbf{Note:} Since p < 0.0001, we report as 0.0000 when rounding to 4 decimal places.

\textbf{Grading:}
\begin{itemize}
\item 0.75 points: Correct \texttt{t.test()} code with correct parameters
\item 0.75 points: Correct p-value (0.0000 or < 0.0001)
\item Acceptable answers: 0.0000, 0, < 0.0001
\end{itemize}

\vspace{1cm}

\subsection*{Question 8: Pooled Standard Deviation (3 points)}

Two independent samples are collected to compare cholesterol levels between two diets:

\begin{itemize}
\item Diet 1: $n_1 = 12$, $s_1 = 18.5$ mg/dL
\item Diet 2: $n_2 = 15$, $s_2 = 21.3$ mg/dL
\end{itemize}

Calculate the pooled standard deviation ($s_p$) using the formula:

$$s_p = \sqrt{\frac{(n_1-1)s_1^2 + (n_2-1)s_2^2}{n_1 + n_2 - 2}}$$

\vspace{0.3cm}

\textbf{Solution:}

\begin{lstlisting}[language=R,basicstyle=\small\ttfamily,frame=single]
# Calculate pooled standard deviation
n1 <- 12
s1 <- 18.5
n2 <- 15
s2 <- 21.3

s_p <- sqrt(((n1-1)*s1^2 + (n2-1)*s2^2) / (n1 + n2 - 2))
s_p
# [1] 20.12421

# Round to 2 decimal places
round(s_p, 2)
# [1] 20.12
\end{lstlisting}

\textbf{Answer:} $s_p$ = \textcolor{red}{\textbf{20.12}} mg/dL

\textbf{Grading:}
\begin{itemize}
\item 1.5 points: Correct formula application (numerator and denominator)
\item 1.5 points: Correct answer (20.12 mg/dL)
\item Partial credit: 19.5-20.5 acceptable for rounding variations
\end{itemize}

\vspace{1cm}

\subsection*{Question 9: Test Statistic Calculation (3 points)}

A researcher compares mean test scores between two teaching methods. The summary statistics are:

\begin{itemize}
\item Method 1: $n_1 = 20$, $\bar{x}_1 = 85.3$, $s_1 = 12.4$
\item Method 2: $n_2 = 18$, $\bar{x}_2 = 78.6$, $s_2 = 14.1$
\end{itemize}

Calculate the test statistic (t-value) for a two-sample t-test (assume equal variances).

\vspace{0.3cm}

\textbf{Solution:}

\begin{lstlisting}[language=R,basicstyle=\small\ttfamily,frame=single]
# Calculate t-statistic from summary statistics
n1 <- 20
xbar1 <- 85.3
s1 <- 12.4

n2 <- 18
xbar2 <- 78.6
s2 <- 14.1

# Step 1: Calculate pooled SD
s_p <- sqrt(((n1-1)*s1^2 + (n2-1)*s2^2) / (n1 + n2 - 2))
s_p
# [1] 13.23003

# Step 2: Calculate standard error
se <- s_p * sqrt(1/n1 + 1/n2)
se
# [1] 4.298345

# Step 3: Calculate t-statistic
t_stat <- (xbar1 - xbar2) / se
t_stat
# [1] 1.558847

# Round to 2 decimal places
round(t_stat, 2)
# [1] 1.56
\end{lstlisting}

\textbf{Answer:} t = \textcolor{red}{\textbf{1.56}}

\textbf{Interpretation:} The t-statistic of 1.56 with df = 36 corresponds to a p-value of approximately 0.13 (two-sided), which is greater than 0.05. Therefore, we would fail to reject the null hypothesis and conclude there is insufficient evidence of a significant difference between the two teaching methods.

\textbf{Grading:}
\begin{itemize}
\item 1.5 points: Correct calculation (pooled SD, SE, then t)
\item 1.5 points: Correct answer (1.56)
\item Partial credit: 1.50-1.60 acceptable for rounding variations
\end{itemize}

\vspace{1cm}

\subsection*{Question 10: Confidence Interval for Mean Difference (3 points)}

Using the data from Question 9 (Teaching Methods), construct a 95\% confidence interval for the mean difference in test scores ($\mu_1 - \mu_2$).

\vspace{0.3cm}

\textbf{Solution:}

\textbf{Method 1: Using summary statistics (manual calculation):}

\begin{lstlisting}[language=R,basicstyle=\small\ttfamily,frame=single]
# From Question 9
n1 <- 20
xbar1 <- 85.3
n2 <- 18
xbar2 <- 78.6
s_p <- 13.23003  # From Q9
se <- 4.298345   # From Q9

# Calculate 95% CI for mean difference
mean_diff <- xbar1 - xbar2
df <- n1 + n2 - 2
t_crit <- qt(0.975, df)

ci_lower <- mean_diff - t_crit * se
ci_upper <- mean_diff + se * t_crit

cat("95% CI: (", round(ci_lower, 2), ",", round(ci_upper, 2), ")\n")
# 95% CI: ( -2.02 , 15.42 )
\end{lstlisting}

\textbf{Method 2: More compact calculation:}

\begin{lstlisting}[language=R,basicstyle=\small\ttfamily,frame=single]
# All in one
mean_diff <- 85.3 - 78.6  # 6.7
t_crit <- qt(0.975, 36)   # 2.028
margin <- 2.028 * 4.298   # 8.72
ci_lower <- 6.7 - 8.72    # -2.02
ci_upper <- 6.7 + 8.72    # 15.42
\end{lstlisting}

\textbf{Answer:} 95\% CI = \textcolor{red}{\textbf{(-2.02, 15.42)}}

\textbf{Interpretation:} We are 95\% confident that the true mean difference in test scores (Method 1 $-$ Method 2) is between $-2.02$ and $15.42$ points. Since the interval \textit{includes zero}, we cannot conclude there is a significant difference between the two teaching methods at $\alpha = 0.05$. The difference could be negative (Method 2 better), zero (no difference), or positive (Method 1 better).

\textbf{Grading:}
\begin{itemize}
\item 1.5 points: Correct CI calculation (formula and computation)
\item 1 point: Correct CI values ($-2.02$, $15.42$)
\item 0.5 points: Correct interpretation (no significant difference, includes 0)
\item Partial credit: CI within $\pm 0.2$ of correct values acceptable
\end{itemize}

\vspace{1cm}

\subsection*{Question 11: Power Calculation (3 points)}

A researcher is planning a one-sample study to test whether a new intervention changes blood glucose levels. Based on prior research:

\begin{itemize}
\item Expected sample size: $n = 30$
\item Expected mean difference to detect: $\delta = 12$ mg/dL
\item Population standard deviation: $\sigma = 18$ mg/dL
\item Significance level: $\alpha = 0.05$ (two-sided)
\end{itemize}

Calculate the statistical power of this study.

\vspace{0.3cm}

\textbf{Solution:}

\begin{lstlisting}[language=R,basicstyle=\small\ttfamily,frame=single]
# Calculate power given sample size and effect size
n <- 30
delta <- 12
sd <- 18
alpha <- 0.05

# Use power.t.test with all parameters specified except power
result <- power.t.test(n = n,
                       delta = delta,
                       sd = sd,
                       sig.level = alpha,
                       type = "one.sample",
                       alternative = "two.sided")

result$power
# [1] 0.9416416

# Round to 2 decimal places
round(result$power, 2)
# [1] 0.94
\end{lstlisting}

\textbf{Answer:} Power = \textcolor{red}{\textbf{0.94}}

\textbf{Interpretation:} With a sample size of 30, this study has 94\% power to detect a mean difference of 12 mg/dL (assuming SD = 18) at the 0.05 significance level. This is excellent power (well above the conventional 80\% threshold), meaning the study has a high probability of correctly rejecting the null hypothesis if the true effect is indeed 12 mg/dL.

\textbf{Grading:}
\begin{itemize}
\item 1.5 points: Correct R code with \texttt{power.t.test()}
\item 1.5 points: Correct answer (0.94)
\item Partial credit: 0.93-0.95 acceptable for rounding variations
\item Common error: Students may report as 94 (percentage) instead of 0.94 - deduct 0.5 points
\end{itemize}

\newpage

\section*{Summary Statistics (for reference)}

\subsection*{Question 6 Verification:}
\begin{verbatim}
Parameters:
  Effect size (delta): 10 mmHg
  Standard deviation: 20 mmHg
  Significance level: 0.05
  Desired power: 0.80

Calculation:
  Exact n = 33.37
  Rounded up = 34 subjects
\end{verbatim}

\subsection*{Question 7 Verification:}
\textbf{Group Means and SDs:}
\begin{verbatim}
Program A: mean = 4.18 kg, SD = 0.259 kg
Program B: mean = 6.08 kg, SD = 0.239 kg
Program C: mean = 7.88 kg, SD = 0.259 kg
\end{verbatim}

\textbf{ANOVA Results:}
\begin{verbatim}
F-statistic: 268.85
df1 (between): 2
df2 (within): 12
p-value: 1.08 × 10^-10 (highly significant)
\end{verbatim}

\textbf{t-Test Results (A vs C):}
\begin{verbatim}
t-statistic: -22.601
df: 8
p-value: 4.457 × 10^-9
Mean difference: 3.70 kg (Program C > Program A)
95% CI: (3.38, 4.02) kg
\end{verbatim}

\textbf{Conclusion:} All three programs differ significantly from each other. Program C is most effective (7.88 kg), followed by Program B (6.08 kg), then Program A (4.18 kg).

\vspace{1cm}

\begin{tcolorbox}[colback=green!5,colframe=green!50!black,boxrule=1pt,arc=0pt]
\textbf{VERIFICATION COMPLETE}

All numerical answers verified using R code in \texttt{verify\_assignment\_answers.R}.

\textbf{MC Answer Key:} 1-B, 2-C, 3-C, 4-C, 5-B

\textbf{Numerical Answer Key:}
\begin{itemize}
\item Question 6: 34 subjects
\item Question 7A: F = 268.85
\item Question 7B: p = 0.0000
\item Question 8: $s_p$ = 20.12 mg/dL
\item Question 9: t = 1.56
\item Question 10: 95\% CI = ($-2.02$, $15.42$)
\item Question 11: Power = 0.94
\end{itemize}

Following CLAUDE.md mandate: NO mental math, ALL values computed with R.
\end{tcolorbox}

\end{document}
